\section{Open Questions and Candidate Next-Step Probes}

Five open questions the paper does not resolve, each with a concrete free-endpoint next-step probe.

\subsection*{Q1. Combined-modality radiation + immunotherapy}
\textbf{Question.} Does the NF-$\kappa$B / COX-2 / PGE\textsubscript{2} axis characterized here in single-agent low-LET X-irradiation of N/TERT-1 3D skin extend to combined radiation-plus-immunotherapy schedules, where PGE\textsubscript{2} is a known immunosuppressive brake on tumor-infiltrating T cells and could shift skin-toxicity risk?

\textbf{Basis.} The paper establishes 2~Gy X-rays alone drive NF-$\kappa$B nuclear translocation, COX-2 induction, and a 6.5$\times$ PGE\textsubscript{2} rise by 72~h in irradiated skin equivalents. Combined-modality radiation + checkpoint-blockade regimens (anti-PD-1, anti-CTLA-4) are now standard of care and have distinct skin-toxicity profiles versus radiation alone. PGE\textsubscript{2} is a well-characterized immunomodulator with dual roles (pro-inflammatory in skin, immunosuppressive in tumor bed). None of this is exercised by the paper, which is single-agent radiation only.

\textbf{Next steps.} Three-arm N/TERT-1 organotypic study: (a) 2~Gy X-rays alone; (b) 2~Gy plus anti-PD-1 or an IFN-$\gamma$ cocktail simulating checkpoint-mediated T-cell activation; (c) 2~Gy plus sc-236 plus IFN-$\gamma$ to test whether COX-2 blockade rescues combined-modality toxicity. Readouts: NF-$\kappa$B nuclear-translocation timecourse by immunofluorescence, COX-2 mRNA/protein by qPCR/Western, PGE\textsubscript{2} by ELISA, K1/FLG differentiation markers, plus IL-6/IL-8 as immune-crosstalk readouts.

\subsection*{Q2. Single-cell NF-$\kappa$B pulsing reconciliation}
\textbf{Question.} How does the population-averaged p-p65 signal reported here (single blot per condition, $n=2$ densitometry) reconcile with the single-cell NF-$\kappa$B pulsing phenotype documented since Nelson et al.\ 2004, and would the paper's Bay 11-7085 dose-response look different if measured on single-cell nuclear-translocation kinetics rather than bulk phospho-p65 by Western?

\textbf{Basis.} Fig 5C reports a monotonic bulk Bay dose-response. Bulk phospho-p65 averages over cells that in single-cell live-imaging studies show digital, oscillatory nuclear translocation with substantial cell-to-cell heterogeneity. A monotonic bulk response can arise either from graded per-cell dampening or from a binary responder-fraction shift, with very different \emph{in vivo} implications.

\textbf{Next steps.} Repeat Fig 5C in N/TERT-1 cells stably expressing RelA-GFP. Live-image nuclear/cytoplasmic p65 ratio for 8~h post-2~Gy at each Bay dose ($n > 100$ cells per condition). Extract per-cell peak amplitude, pulse frequency, and responder fraction. Fit a stochastic IKK/I$\kappa$B/NF-$\kappa$B module (Hoffmann 2002 topology) to the single-cell traces and predict bulk COX-2 mRNA at each dose.

\subsection*{Q3. Radiotherapy fractionation}
\textbf{Question.} Does the 2~Gy single-fraction protocol used here predict skin toxicity under clinically relevant fractionation schedules (2~Gy $\times$ 30 conventional, 3~Gy $\times$ 15 hypofractionated, 8~Gy $\times$ 1 SBRT), and does sc-236 rescue efficacy depend on fractionation?

\textbf{Basis.} Clinical radiotherapy for breast, head-and-neck, and cutaneous indications uses fractionated regimens spanning 2--8~Gy per fraction and 5--30 fractions total. The paper's 2~Gy single-fraction result does not touch the dose-per-fraction or inter-fraction repair biology that governs clinical skin toxicity. NF-$\kappa$B has known adaptive/refractory kinetics on repeated stimulation.

\textbf{Next steps.} Three fractionation arms in the N/TERT-1 organotypic protocol: 2~Gy $\times$ 5 daily, 3~Gy $\times$ 3, 8~Gy $\times$ 1. Measure COX-2 mRNA and PGE\textsubscript{2} at 24~h post-final fraction. Add a sc-236 continuous-dosing arm for each schedule. Compare to LQ-model EQD2 predictions and cross-reference clinical RTOG skin-toxicity grading data from TCIA/NCTN Data Archive.

\subsection*{Q4. Feedback-loop topology sensitivity}
\textbf{Question.} How sensitive is the paper's implicit signaling-topology assumption (linear NF-$\kappa$B $\to$ COX-2 $\to$ PGE\textsubscript{2} $\to$ differentiation) to alternative topologies with feedback loops --- specifically, does adding a PGE\textsubscript{2}-to-NF-$\kappa$B positive-feedback arc via EP2/EP4 receptors change the predicted response to sc-236 versus Bay?

\textbf{Basis.} The paper implicitly assumes a feed-forward cascade with sc-236 acting downstream (COX-2 block) and Bay acting upstream (NF-$\kappa$B block). Both rescue keratinocyte differentiation, consistent with a linear model. But the inflammation literature documents PGE\textsubscript{2} signaling back to NF-$\kappa$B via EP2/EP4 receptors and cAMP/PKA, closing a positive-feedback loop that would give sc-236 and Bay quantitatively different upstream effects.

\textbf{Next steps.} Minimal 4-species ODE model in Python: NF-$\kappa$B active fraction, COX-2 protein, PGE\textsubscript{2} extracellular, K1/FLG index. Parameterize from Fig 1 and Fig 7A timecourses. Fit two topologies: (T1) feed-forward only; (T2) T1 plus a PGE\textsubscript{2}$\to$NF-$\kappa$B positive-feedback arc with EP2/EP4 rate constants from published inflammation-signaling literature. Compare predicted sc-236 and Bay dose-response shapes to Fig 5B/5C. Report $\Delta$AIC between topologies and predicted p-p65 fold-difference at each Bay dose.

\subsection*{Q5. ML-driven skin-toxicity prediction}
\textbf{Question.} Can an ML model trained on paired skin-imaging and gene-expression endpoints (H\&E histology, IHC panels for K1/FLG/COX-2, and bulk RNA-seq of irradiated organotypic skin) predict radiation-induced skin toxicity grades better than either modality alone, and would the model recover the NF-$\kappa$B / COX-2 axis identified biologically?

\textbf{Basis.} The paper connects a molecular readout (COX-2 mRNA, p-p65, PGE\textsubscript{2}) to a phenotypic readout (K1/FLG immunofluorescence, cornified-layer thickness) on the same organotypic samples. This is exactly the paired imaging plus expression structure current multimodal ML architectures (CLIP-style contrastive learning, graph-neural histology plus expression models) are designed to exploit. Clinical RTOG skin-toxicity grading currently relies on qualitative visual assessment.

\textbf{Next steps.} Assemble a training set from public repositories: TCIA skin-toxicity IHC panels; GEO/ArrayExpress irradiated-keratinocyte bulk RNA-seq. Train a two-tower contrastive model (histolab tissue tiling, scanpy expression preprocessing, PyTorch multimodal head). Evaluate ROC on RTOG grade $\geq 2$ on held-out samples. Use SHAP / Integrated Gradients to verify the model attends to COX-2, PTGS2, NFKB1, RELA gene features and to cornified-layer thickness in H\&E --- positive controls that it recovers the paper's biology.
